\chapter{Materials and Methods}
\label{ch:methods}

\section{Interfacial polymerisation}

The polyamide layer is formed by reacting \emph{m}-phenylenediamine (MPD) with
trimesoyl chloride (TMC). The condensation reaction releases hydrogen chloride
(\ce{HCl}):

\begin{equation}
  \ce{R-NH2 + ClOC-R2 -> R-NH-CO-R2 + HCl}
  \label{eq:ip}
\end{equation}

where R and R2 denote the aromatic backbones of MPD and TMC respectively.
Chemical formulae use \texttt{mhchem}, e.g.\ sodium chloride \ce{NaCl} and
the sulfate ion \ce{SO4^2-}.

\section{Performance metrics}

Pure water permeability is calculated as

\begin{equation}
  A = \frac{J_w}{\Delta P} = \frac{V}{A_m \, t \, \Delta P},
  \label{eq:pwp}
\end{equation}

\noindent where $J_w$ is the water flux, $\Delta P$ the applied pressure,
$V$ the permeate volume, $A_m$ the membrane area and $t$ time. Observed salt
rejection is

\begin{equation}
  R = \left(1 - \frac{C_p}{C_f}\right) \times 100\%,
  \label{eq:rejection}
\end{equation}

\noindent with $C_p$ and $C_f$ the permeate and feed concentrations.

Units are typeset with \texttt{siunitx}: a feed pressure of
\SI{15}{\bar} and a conductivity of \SI{2000}{\micro\siemens\per\centi\metre}.

\section{Example table}

\begin{table}[htbp]
  \centering
  \caption{Example membrane performance data.}
  \label{tab:performance}
  \begin{tabular}{lcc}
    \toprule
    Sample & PWP (\si{\litre\per\square\metre\per\hour\per\bar}) & Rejection (\%) \\
    \midrule
    Virgin       & 3.8 & 98.7 \\
    End-of-life   & 9.1 & 41.2 \\
    Regenerated  & 4.2 & 96.5 \\
    \bottomrule
  \end{tabular}
\end{table}

Refer to tables and equations by number: Table~\ref{tab:performance} and Equation~\eqref{eq:rejection}.
